\section{Divisor multiplication and degree with the full lattice fibre}
\label{sec:full-lattice-divisors}

The source construction is Alain Connes and Caterina Consani,
\href{https://arxiv.org/abs/1502.05585v2}{\emph{Absolute algebra and
Segal's $\Gamma$-sets}}, original \texttt{F1\_corr.tex},
\S5.2: Proposition \texttt{propspecz}, equations \texttt{propsmash1},
\texttt{propsmash2}, and Proposition \texttt{proppic}, source lines
938--1053. Their divisor modules and exceptional multiplication rule
are the source results. Here we calculate the rule for the full
lattice carrier and replace the global-section test by an intrinsic
test that retains every labelled zero. The split carrier is from
The Clankers, \emph{Split Support Geometry}, v11,
\href{https://github.com/KokunoYumeto/zeta-function-research-reader/blob/a2851f9eb352e7e4376bc629de86fa4ed9c1c434/workbenches/splitzero-tandem/foundations/split-support-geometry-v11/split_support_geometry_arithmetic_curve_v11.tex#L1226}{Definition \texttt{def:lattice-split}}.

\subsection{The original divisor and its entire support lift}
Fix a nontrivial bounded distributive lattice $L$ and write
$z_\ell=(0,\ell)$, $\tau=z_{0_L}$, $e=z_{1_L}$, and
$\widehat x=(x,1_L)$ for $x\ne0$. Retain the original divisor
\begin{equation}\tag{SDM1}
 D=\sum_p n_p[p]+a[\infty],\qquad n_p\in\mathbb Z,
 \quad n_p=0\text{ for all but finitely many }p,\quad a\in\mathbb R.
\end{equation}
The underlying space is the source's
$X=\overline{\operatorname{Spec}\mathbb Z}$, whose nonempty opens
are complements of finite sets of closed primes, possibly including
$\infty$. This section constructs sheaves on that specified space;
it does not identify it with $\operatorname{Spec}G_L(\mathbb Z)$.
For nonempty $U$ define
\begin{equation}\tag{SDM2}
 \begin{split}
 R_U&=\{x\in\mathbb Q:v_p(x)\ge0\text{ for finite }p\in U\},\\
 I_D(U)&=\{x\in\mathbb Q:v_p(x)+n_p\ge0\text{ for finite }p\in U\},\\
 G_L(I_D(U))&=\{(0,\ell):\ell\in L\}
           \cup\{(x,1_L):x\in I_D(U)\}.
 \end{split}
\end{equation}
Here $v_p(0)=+\infty$. Addition joins the supports, and the scalar
action of $G_L(R_U)$ multiplies amplitudes and meets supports.
For a finite pointed set $K$, with $K^\circ=K\setminus\{*\}$,
put
\begin{equation}\tag{SDM3}
 \mathcal M_L(D)(U)(K)=
 \left\{(x_j,\ell_j)_{j\in K^\circ}\in G_L(I_D(U))^{K^\circ}:
 \ \infty\in U\Longrightarrow\sum_j|x_j|\le e^a\right\}.
\end{equation}
The empty-open value is the singleton. A pointed map sums the
amplitudes and joins the supports over each nonbasepoint fibre;
an empty fibre has value $\tau$. Restriction includes the rational
amplitudes and leaves every support unchanged. Write
$\mathcal A_L=\mathcal M_L(0)$. Also define
$\mathcal M_L^{<}(D)$ by replacing $\le e^a$ by $<e^a$ on
opens containing $\infty$, with no change elsewhere.

\begin{theorem}\label{SDM:sheaf}
These are sheaves of $\mathcal A_L$-modules. The amplitude-zero
subsheaf is $HL$ on every nonempty open, for both inequalities.
Their stalks are
\begin{equation}\tag{SDM4}
 \begin{aligned}
 \mathcal M_L(D)_\eta&=HG_L(\mathbb Q),\\
 \mathcal M_L(D)_p&=HG_L(p^{-n_p}\mathbb Z_{(p)}),\\
 \mathcal M_L(D)_\infty&=B_{e^a}^L(\mathbb Q).
 \end{aligned}
\end{equation}
For the strict sheaf the last stalk is $B_{<e^a}^L(\mathbb Q)$.
The notation $H$ on an additive module denotes its finite-tuple
$\Gamma$-set; it does not assert that $I_D(U)$ is a ring.
\end{theorem}
\begin{proof}
The valuation inequalities make $I_D(U)$ an $R_U$-module. Addition
and scalar action on the displayed carrier satisfy their axioms
by distributivity in $L$ and in $\mathbb Q$. A nonzero product has
both input supports equal to $1_L$, as required. The triangle
inequality proves closure under pointed maps, including maps that
discard input coordinates. For the external module action the
coordinates are $(r_i x_j,\lambda_i\wedge\ell_j)$; its bound is
$\sum_{i,j}|r_i x_j|=(\sum_i|r_i|)(\sum_j|x_j|)\le e^a$,
and is strict when the second bound is strict. The valuation
inequalities add. This proves the module identities coordinatewise.

Every two nonempty opens intersect. Compatible local sections
therefore have exactly the same rational amplitudes and lattice
labels. All valuation conditions required on the union hold on a
covering open; if the union contains infinity, one such open supplies
the required bound. This proves existence and uniqueness of gluing.
Every zero-amplitude tuple satisfies either positive bound, proving
the subsheaf assertion. At a finite prime $p$, all other primes $q$
where one of the finitely many inequalities $v_q(x_j)+n_q\ge0$
fails can be excluded, together with infinity. This is a finite
set: it includes denominator primes and may also include primes with
negative divisor coefficient. At infinity every such finite prime
can be excluded, while the original bound remains. At the generic
point every such finite prime and infinity can be excluded.
Taking these unions gives (SDM4) without changing any label.
\end{proof}

\subsection{The exact tensor law, including its boundary}
Let $D'=\sum_p n'_p[p]+a'[\infty]$, and set
$\lambda=e^a$, $\mu=e^{a'}$. Multiplication defines
\begin{equation}\tag{SDM5}
 m_{D,D'}:\mathcal M_L(D)\wedge_{\mathcal A_L}\mathcal M_L(D')
 \longrightarrow\mathcal M_L(D+D'),\qquad
 (x_i,\ell_i),(y_j,t_j)\longmapsto(x_i y_j,\ell_i\wedge t_j).
\end{equation}
The relative smash product is the sheafification of the module
coequalizer for the Day smash product. The displayed product is
balanced and respects all pointed maps by the two distributive laws.

\begin{theorem}\label{SDM:tensor}
The map (SDM5) is injective as a map of sheaves and identifies its
source with
\begin{equation}\tag{SDM6}
 \begin{cases}
 \mathcal M_L^{<}(D+D'),&\lambda,\mu\notin\mathbb Q
                             \text{ and }\lambda\mu\in\mathbb Q,\\
 \mathcal M_L(D+D'),&\text{otherwise}.
 \end{cases}
\end{equation}
In particular its image contains the whole $HL$ in every case.
The exceptional missing stalk tuples are exactly those at infinity
whose rational amplitudes satisfy $\sum_j|x_j|=\lambda\mu$;
every permissible choice of support on their zero coordinates is
missing with them. The finite stalks have no such defect.
\end{theorem}
\begin{proof}
We give the module calculation, rather than infer it from level one.
At a finite prime, multiplication by the rational $p^{-n_p}$ gives
an isomorphism from $HG_L(\mathbb Z_{(p)})$ to the second object
in (SDM4). The map multiplies every amplitude and fixes every label;
its inverse multiplies by $p^{n_p}$. The same holds for $D'$.
The unit identity $A\wedge_A A\simeq A$ now proves the claimed
isomorphism on that stalk. At the generic point both modules are
the coefficient module itself.

At infinity put $A=B_1^L(\mathbb Q)$. For each rational $r>0$,
\begin{equation}\tag{SDM7}
 s_r:A\xrightarrow{\sim}B_r^L(\mathbb Q),\qquad
 (x_j,\ell_j)_j\longmapsto(r x_j,\ell_j)_j
\end{equation}
is an $A$-module isomorphism, with inverse $s_{1/r}$ on the indicated
codomain. For $0<r\le t$ the inclusion $B_r^L\subset B_t^L$
corresponds under these isomorphisms to multiplication by
$\widehat{r/t}$ in $A(1_+)$. Thus the multiplication map induces
$B_r^L\wedge_A B_s^L\simeq B_{rs}^L$ and the transition maps are
exactly the coordinatewise inclusions, including all zero labels.
For completeness, the inverse of the balanced-unit multiplication
is $A\simeq A\wedge\mathbb S\to A\wedge A\to A\wedge_A A$,
using the unit $\mathbb S\to A$ in the middle map. One composite
is the identity by the unit law; the other is the identity because
the balanced relation transfers the first $A$-factor to the second.
A tuple $(z_{\ell_j})_j$ is represented by that tuple paired with
$\widehat1$ and is recovered with supports $\ell_j\wedge1_L=\ell_j$.

For any real $c>0$ the full bounded module is the filtered union
\begin{equation}\tag{SDM8}
 B_c^L(\mathbb Q)=\underset{r\in\mathbb Q,\ 0<r\le c}{\operatorname{colim}}
 B_r^L(\mathbb Q).
\end{equation}
Indeed each tuple has rational total absolute amplitude. If that
total is positive it may be used as $r$; if it is zero any rational
$0<r\le c$ works, retaining its arbitrary lattice labels. The
relative smash product preserves colimits in each variable: the
Day smash is a left adjoint in each variable, and the defining
balanced coequalizer commutes with colimits. Hence the product at
infinity is the union of $B_{rs}^L$ over positive rational
$r\le\lambda$, $s\le\mu$. A tuple of rational total amplitude
$t<\lambda\mu$ belongs to this union by choosing $r,s$ sufficiently
close to their upper endpoints. The same is true for $t=0$.
The endpoint $t=\lambda\mu$ is attained precisely when both
endpoints $\lambda,\mu$ are rational: since $r/\lambda,s/\mu\le1$,
their product is one only when both equal one. If exactly one
endpoint is irrational the product is irrational, so no rational
tuple total equals it. If both are irrational and their product is
irrational, the same observation applies. This proves (SDM6) at
infinity. All these are injections into the same tuple module.

Stalks commute with the tensor construction used here. To verify
this directly, the stalk is a filtered colimit over neighborhoods;
each Day representative uses finite pointed sets and finitely many
sections, so all its entries occur on a common neighborhood.
Every defining relation between representatives also occurs on a
common neighborhood. The same finite-representative argument applies
to the balanced coequalizer. Sheafification leaves stalks unchanged.
Consequently the stalk computations prove the sheaf statement.
\end{proof}

\subsection{Recovering degree without deleting labelled zeros}
For a pointed $\Gamma$-set $M$, call $x\in M(1_+)$ additively
idempotent when there exists $w\in M(2_+)$ with
\begin{equation}\tag{SDM9}
 M(\pi_1)w=x=M(\pi_2)w,\qquad M(\nabla)w=x,
\end{equation}
where $\pi_1,\pi_2$ are the two pointed coordinate projections
and $\nabla$ is the fold. This property uses only the given
$\Gamma$-structure and is preserved by isomorphisms.

\begin{lemma}\label{SDM:idempotents}
For every nonempty open in (SDM3), its additively idempotent
level-one elements are exactly $\{z_\ell:\ell\in L\}$, for either
inequality. The same holds for the image sheaf in (SDM6).
\end{lemma}
\begin{proof}
The two projections force $w=(x,x)$. If $x=(u,\ell)$, the fold
equation says $2u=u$ and $\ell\vee\ell=\ell$. In $\mathbb Q$
the first equality is equivalent to $u=0$. Conversely $(z_\ell,z_\ell)$
always satisfies the valuation conditions and either positive bound,
and its fold is $z_\ell$. The image description (SDM6) supplies the
last assertion.
\end{proof}

Keep all the coefficients of (SDM1), and define
\begin{equation}\tag{SDM10}
 c_D=\prod_p p^{-n_p}>0,\qquad
 \deg D=a+\sum_p n_p\log p=a-\log c_D.
\end{equation}
These formulas do not replace $D$ by another divisor. The actual
global amplitude module $I_D(X)$ is $c_D\mathbb Z$: for a rational
$x$, all $v_p(x/c_D)\ge0$ hold exactly when $x/c_D\in\mathbb Z$.
For $D_b=b[\infty]$, set
$\mathcal T_b=\mathcal M_L(D)\wedge_{\mathcal A_L}\mathcal M_L(D_b)$.
Define the following invariant of an abstract
$\mathcal A_L$-module sheaf of the displayed form:
\begin{equation}\tag{SDM11}
 \alpha_L(\mathcal M_L(D))=
 \inf\left\{b\in\mathbb R:
 \mathcal T_b(X)(1_+)
 \text{ contains a non-idempotent element}\right\}.
\end{equation}
\begin{theorem}\label{SDM:degree}
One has the exact formula and classification
\begin{equation}\tag{SDM12}
 \alpha_L(\mathcal M_L(D))=\log c_D-a=-\deg D,
 \qquad
 \mathcal M_L(D)\simeq\mathcal M_L(D')
 \Longleftrightarrow D'-D\text{ is principal}.
\end{equation}
The invariant in (SDM11) uses the complete support sheaf and no
quotient of its zero fibre.
\end{theorem}
\begin{proof}
By (SDM6), a non-idempotent global section is a nonzero element of
$c_D\mathbb Z$ with $|x|\le e^{a+b}$, with strict inequality in
the exceptional case. Its least possible absolute value is exactly
$c_D$. Thus it does not exist for $b<\log c_D-a$ and does exist
for every $b>\log c_D-a$. At equality it exists precisely when
the tensor law is not exceptional, a distinction that leaves the
infimum unchanged. More precisely, the set in (SDM11) is
$[-\deg D,\infty)$ when $e^a\in\mathbb Q$ and
$(-\deg D,\infty)$ otherwise: at equality the product of the
two radius bounds is $c_D\in\mathbb Q$, and $e^b=c_D/e^a$ is
rational exactly when $e^a$ is rational. This proves the first
formula with the endpoint fully accounted for. An isomorphism of module sheaves induces an
isomorphism after tensoring with $\mathcal M_L(D_b)$; it preserves
(SDM9), and therefore preserves this infimum. It follows that
isomorphic sheaves have $\deg D=\deg D'$.

Put $q=\prod_p p^{n'_p-n_p}\in\mathbb Q_{>0}$. Equality of degrees
gives $a'-a=-\log q$, while $v_p(q)=n'_p-n_p$ for every $p$.
Thus $D'=D+(q)$ with the original principal-divisor convention
$(q)=\sum_pv_p(q)[p]-\log|q|[\infty]$. Conversely, for any
$q\in\mathbb Q^\times$ and $D'=D+(q)$, the explicit map
\begin{equation}\tag{SDM13}
 \mathcal M_L(D)(U)(K)\longrightarrow\mathcal M_L(D')(U)(K),
 \qquad(x_j,\ell_j)_j\longmapsto(q^{-1}x_j,\ell_j)_j
\end{equation}
is an isomorphism. Indeed
$v_p(q^{-1}x_j)+n_p+v_p(q)=v_p(x_j)+n_p$, and the new
archimedean bound is $e^a/|q|$. Its inverse multiplies by $q$.
It commutes with all restrictions, pointed maps and module actions,
and fixes every $z_\ell$. This proves both implications.
\end{proof}

For nontrivial $L$, replacing ``non-idempotent'' in (SDM11) by
``different from the pointed zero'' would give $-\infty$, since
$e\ne\tau$ is always present. This is the precise defect of
transporting the source's nonzero-section test without its lattice
fibre. The intrinsic test proves (SDM12) while preserving that fibre.
The tensor calculation concerns the original sum-of-absolute-values
divisor sheaves. It does not identify them with the later all-subset
modules of \emph{Riemann--Roch for the ring $\mathbb Z$}; their exact
comparison is (HG38)--(HG41). No weight or cohomological purity
statement is inferred from these module calculations.
